\documentclass[12pt,a4paper]{ctexart}
\usepackage[top=2.5cm,bottom=2.5cm,left=2.5cm,right=2.5cm]{geometry}
\usepackage{amsmath,amssymb,amsfonts}
\usepackage{graphicx}
\usepackage{tikz}
\usetikzlibrary{shapes,arrows,positioning,calc,decorations.markings,matrix}
\usepackage{xcolor}
\usepackage{enumitem}
\usepackage{booktabs}
\usepackage{hyperref}
\usepackage{physics}
\usepackage{float}
\usepackage{caption}
\usepackage{framed}
\usepackage{tcolorbox}
\tcbuselibrary{skins,breakable}

\hypersetup{
    colorlinks=true,
    linkcolor=blue!70!black,
    citecolor=green!50!black,
    urlcolor=blue!70!black
}

% 定义颜色
\definecolor{tensorblue}{RGB}{70,130,180}
\definecolor{tensorred}{RGB}{205,92,92}
\definecolor{tensorgreen}{RGB}{60,179,113}
\definecolor{tensororange}{RGB}{230,160,50}
\definecolor{tensorpurple}{RGB}{147,112,219}
\definecolor{queencolor}{RGB}{220,50,50}

% TikZ 样式
\tikzset{
    tensor/.style={draw, thick, minimum size=0.6cm},
    cleg/.style={thick, tensorblue},
    pleg/.style={thick, tensorred},
    queen/.style={fill=queencolor!80, text=white, font=\bfseries\large},
    empty/.style={fill=gray!10},
}

\title{\textbf{$N$ 皇后问题的张量网络解法\\——从约束编码到精确计数}}
\author{}
\date{}

\begin{document}
\maketitle
\thispagestyle{empty}

\begin{abstract}
$N$ 皇后问题是组合数学中的经典问题：在 $N \times N$ 的棋盘上放置 $N$ 个皇后，使得任意两个皇后不能互相攻击（不在同一行、同一列、同一对角线上）。本文详细介绍如何将 $N$ 皇后的\textbf{计数问题}（\#$N$-Queens）转化为一个张量网络的收缩问题，并展示从约束编码、张量构造、转移矩阵的 MPO 分解，到精确与近似收缩的完整方法链。通过一个 $N=4$ 的显式计算实例，我们将每一步骤都具体化。最后，我们讨论推广到皇后格点气体模型及其相变行为。
\end{abstract}

\tableofcontents
\newpage

%========================================
\section{引言}
%========================================

\subsection{$N$ 皇后问题}

$N$ 皇后问题的\textbf{判定版本}是：能否在 $N \times N$ 棋盘上放置 $N$ 个互不攻击的皇后？对于 $N \geq 4$（以及 $N=1$），答案总是"能"。

更有趣也更困难的是\textbf{计数版本}（\#$N$-Queens）：这样的摆放方式一共有多少种？下表列出了已知的精确值：

\begin{center}
\renewcommand{\arraystretch}{1.2}
\begin{tabular}{c|cccccccccc}
\hline
$N$ & 1 & 2 & 3 & 4 & 5 & 6 & 7 & 8 & 9 & 10 \\
\hline
$Z_N$ & 1 & 0 & 0 & 2 & 10 & 4 & 40 & 92 & 352 & 724 \\
\hline\hline
$N$ & 11 & 12 & 13 & 14 & 15 & 16 & 17 & 18 & & \\
\hline
$Z_N$ & 2680 & 14200 & 73712 & 365596 & 2279184 & 14772512 & 95815104 & 666090624 & & \\
\hline
\end{tabular}
\end{center}

\textbf{计算复杂性：} \#$N$-Queens 问题属于 \#P-完全类——这意味着它至少和 NP-完全问题一样困难，很可能不存在多项式时间算法。经典的回溯搜索（backtracking）对于 $N \lesssim 27$ 是可行的，但计算量随 $N$ 指数增长。

\subsection{为什么用张量网络？}

张量网络方法提供了一种\textbf{系统化的框架}来处理这类计数问题：

\begin{enumerate}
    \item \textbf{将约束编码为张量}：每个约束（行、列、对角线上至多/恰好一个皇后）都可以精确地表示为一个小张量。
    \item \textbf{组合问题 $\leftrightarrow$ 张量网络收缩}：解的总数等于所有约束张量组成的网络的完全收缩值，形式上与统计物理中的配分函数完全一致。
    \item \textbf{转移矩阵天然具有 MPO 结构}：逐行处理棋盘时，转移矩阵可以分解为矩阵乘积算符（MPO），使得计算结构化。
    \item \textbf{可控的近似}：对于精确求解不可行的大 $N$，可以通过截断键维度（类似 DMRG 中的截断）获得近似结果，误差可控。
\end{enumerate}

本文的目标是：从最基本的问题定义出发，\textbf{逐步构造}这个张量网络，使每个步骤都清晰可循。

%========================================
\section{问题的数学表述}
%========================================

\subsection{二值变量与约束}

在 $N \times N$ 的棋盘上，为每个格点 $(i,j)$（$i$ = 行，$j$ = 列，$1 \leq i,j \leq N$）引入一个\textbf{二值变量}：
\begin{equation}
n_{ij} = \begin{cases} 1, & \text{格点 $(i,j)$ 上有皇后} \\ 0, & \text{格点 $(i,j)$ 上没有皇后} \end{cases}
\end{equation}

$N$ 皇后的合法摆放必须满足以下\textbf{三类约束}：

\subsubsection*{（A）行约束——每行恰好一个皇后}
\begin{equation}
\sum_{j=1}^{N} n_{ij} = 1, \quad \forall\, i = 1, \ldots, N
\end{equation}

\subsubsection*{（B）列约束——每列恰好一个皇后}
\begin{equation}
\sum_{i=1}^{N} n_{ij} = 1, \quad \forall\, j = 1, \ldots, N
\end{equation}

\subsubsection*{（C）对角线约束——每条对角线至多一个皇后}

正方格子上有两族对角线：
\begin{itemize}
    \item \textbf{"$/$" 对角线}（反对角线）：由 $i + j = \text{常数}$ 定义，共 $2N - 1$ 条。
    \item \textbf{"$\backslash$" 对角线}（主对角线方向）：由 $i - j = \text{常数}$ 定义，共 $2N - 1$ 条。
\end{itemize}

每条对角线 $d$ 上的约束为：
\begin{equation}
\sum_{(i,j) \in d} n_{ij} \leq 1
\end{equation}

\subsection{计数问题作为"配分函数"}

$N$ 皇后解的总数可以写成所有构型的加权求和——只不过权重要么为 1（合法），要么为 0（非法）：

\begin{equation}
\boxed{
Z_N = \sum_{\{n_{ij}\}} \prod_{i=1}^{N} \underbrace{\delta\!\left(\sum_j n_{ij} = 1\right)}_{\text{行约束}} \cdot \prod_{j=1}^{N} \underbrace{\delta\!\left(\sum_i n_{ij} = 1\right)}_{\text{列约束}} \cdot \prod_{d \in \mathcal{D}} \underbrace{\delta\!\left(\sum_{(i,j)\in d} n_{ij} \leq 1\right)}_{\text{对角线约束}}
}
\end{equation}

这里 $\delta(\text{条件})$ 是指示函数（条件满足时为 1，否则为 0），$\mathcal{D}$ 是所有对角线的集合。

\textbf{关键观察}：这个表达式的结构和统计力学中的\textbf{配分函数}一模一样！每个约束对应一个"Boltzmann 权重"（取值 0 或 1），总数 $Z_N$ 就是所有"微观态"的权重之和。这不是巧合——这种对应是张量网络方法能够处理组合问题的根本原因。

%========================================
\section{约束的张量网络编码}
%========================================

本节是全文的技术核心：我们将展示如何把上述每种约束\textbf{精确地}编码为一个矩阵乘积态（MPS），从而为构造完整的张量网络铺平道路。

\subsection{"恰好一个"约束的 MPS 表示}

\textbf{任务}：给定 $N$ 个二值变量 $n_1, n_2, \ldots, n_N \in \{0, 1\}$，构造一个函数
\[
f_{\text{exact}}(n_1, \ldots, n_N) = \delta\!\left(\sum_{k=1}^N n_k = 1\right)
\]
并将其表示为 MPS。

\textbf{思路}：从左到右扫描 $n_1, n_2, \ldots, n_N$，用一个\textbf{内部状态} $\alpha$ 记录"到目前为止已经出现了多少个 1"。由于我们只关心总数是否为 1，状态空间只需两个值：

\begin{itemize}
    \item $\alpha = 0$：到目前为止还没有出现过 1（没有皇后）
    \item $\alpha = 1$：已经出现了恰好一个 1（已放一个皇后）
\end{itemize}

如果在 $\alpha = 1$ 的状态下再遇到 $n_k = 1$（第二个皇后），转移权重为 0——直接"杀死"这条路径。

定义 $2 \times 2$ 的\textbf{局部张量}（对每个 $n_k$ 的取值各一个矩阵）：
\begin{equation}
\boxed{
A^{[0]} = \begin{pmatrix} 1 & 0 \\ 0 & 1 \end{pmatrix}, \qquad
A^{[1]} = \begin{pmatrix} 0 & 1 \\ 0 & 0 \end{pmatrix}
}
\end{equation}

\begin{itemize}
    \item $A^{[0]}$：$n_k = 0$（无皇后），状态不变——$0 \to 0$，$1 \to 1$。
    \item $A^{[1]}$：$n_k = 1$（有皇后），只允许 $0 \to 1$（第一个皇后）。从状态 1 出发的所有转移为 0（第二个皇后被禁止）。
\end{itemize}

加上边界向量：
\begin{equation}
\langle L| = (1, 0), \qquad |R\rangle = \begin{pmatrix} 0 \\ 1 \end{pmatrix}
\end{equation}
$\langle L|$：从"0 个皇后"的状态出发。$|R\rangle$：要求最终状态为"恰好 1 个皇后"。

于是：
\begin{equation}
\boxed{
f_{\text{exact}}(n_1, \ldots, n_N) = \langle L| \, A^{[n_1]} A^{[n_2]} \cdots A^{[n_N]} \, |R\rangle
}
\end{equation}

\textbf{键维度 $D = 2$}——和变量个数 $N$ 无关！

\subsubsection*{验证（$N=4$ 为例）}

取 $(n_1, n_2, n_3, n_4) = (0, 1, 0, 0)$（第 2 列放皇后）：
\begin{align}
&\langle L| A^{[0]} A^{[1]} A^{[0]} A^{[0]} |R\rangle \notag \\
&= (1,0) \begin{pmatrix}1&0\\0&1\end{pmatrix} \begin{pmatrix}0&1\\0&0\end{pmatrix} \begin{pmatrix}1&0\\0&1\end{pmatrix} \begin{pmatrix}1&0\\0&1\end{pmatrix} \begin{pmatrix}0\\1\end{pmatrix} \notag \\
&= (1,0) \begin{pmatrix}0&1\\0&0\end{pmatrix} \begin{pmatrix}0\\1\end{pmatrix} = (1,0) \begin{pmatrix}1\\0\end{pmatrix} = 1 \quad \checkmark
\end{align}

取 $(n_1, n_2, n_3, n_4) = (1, 0, 1, 0)$（第 1 和第 3 列各放一个，违规！）：
\begin{align}
&\langle L| A^{[1]} A^{[0]} A^{[1]} A^{[0]} |R\rangle \notag \\
&= (1,0) \begin{pmatrix}0&1\\0&0\end{pmatrix} \begin{pmatrix}1&0\\0&1\end{pmatrix} \begin{pmatrix}0&1\\0&0\end{pmatrix} \begin{pmatrix}1&0\\0&1\end{pmatrix} \begin{pmatrix}0\\1\end{pmatrix} \notag \\
&= (0,1) \begin{pmatrix}0&1\\0&0\end{pmatrix} \begin{pmatrix}0\\1\end{pmatrix} = (0,1) \begin{pmatrix}1\\0\end{pmatrix} = 0 \quad \checkmark
\end{align}

取 $(0,0,0,0)$（不放皇后）：
\[
(1,0) \begin{pmatrix}1&0\\0&1\end{pmatrix}^4 \begin{pmatrix}0\\1\end{pmatrix} = (1,0)\begin{pmatrix}0\\1\end{pmatrix} = 0 \quad \checkmark
\]

\subsection{"至多一个"约束的 MPS 表示}

对角线约束要求每条对角线上\textbf{至多}一个皇后（允许零个）。编码方式几乎完全一样，只需修改右边界向量：

\begin{equation}
\boxed{
g_{\text{atmost}}(n_1, \ldots, n_N) = \langle L| \, A^{[n_1]} \cdots A^{[n_N]} \, |R'\rangle, \qquad |R'\rangle = \begin{pmatrix} 1 \\ 1 \end{pmatrix}
}
\end{equation}

$|R'\rangle = (1,1)^T$ 意味着最终状态为 0 或 1 都可以接受——即至多一个皇后。

键维度仍为 $D = 2$。

\subsection{约束的图形表示}

用张量网络图来表示，每个约束是一条"链"，穿过它所覆盖的所有格点：

\begin{figure}[H]
\centering
\begin{tikzpicture}[scale=0.9]
    % Row constraint
    \node[font=\small\bfseries] at (-1.5, 2) {行约束：};
    \node[tensor, fill=yellow!20, minimum size=0.6cm] (L1) at (0,2) {$\langle L|$};
    \foreach \j in {1,2,3,4} {
        \node[tensor, fill=blue!15, minimum size=0.8cm] (R\j) at (\j*1.5, 2) {$A^{[n_{i\j}]}$};
    }
    \node[tensor, fill=yellow!20, minimum size=0.6cm] (R5) at (7.5,2) {$|R\rangle$};
    \draw[cleg] (L1) -- (R1);
    \draw[cleg] (R1) -- (R2);
    \draw[cleg] (R2) -- (R3);
    \draw[cleg] (R3) -- (R4);
    \draw[cleg] (R4) -- (R5);

    % Column constraint
    \node[font=\small\bfseries] at (-1.5, 0) {列约束：};
    \node[tensor, fill=yellow!20, minimum size=0.6cm] (L2) at (0,0) {$\langle L|$};
    \foreach \i in {1,2,3,4} {
        \node[tensor, fill=green!15, minimum size=0.8cm] (C\i) at (\i*1.5, 0) {$A^{[n_{\i j}]}$};
    }
    \node[tensor, fill=yellow!20, minimum size=0.6cm] (C5) at (7.5,0) {$|R\rangle$};
    \draw[tensorgreen, thick] (L2) -- (C1);
    \draw[tensorgreen, thick] (C1) -- (C2);
    \draw[tensorgreen, thick] (C2) -- (C3);
    \draw[tensorgreen, thick] (C3) -- (C4);
    \draw[tensorgreen, thick] (C4) -- (C5);

    % Diagonal constraint
    \node[font=\small\bfseries] at (-1.5, -2) {对角线约束：};
    \node[tensor, fill=yellow!20, minimum size=0.6cm] (L3) at (0,-2) {$\langle L|$};
    \foreach \k [count=\ki] in {1,2,3} {
        \node[tensor, fill=red!10, minimum size=0.8cm] (D\ki) at (\ki*1.5, -2) {$B^{[n_d]}$};
    }
    \node[tensor, fill=yellow!20, minimum size=0.6cm] (D4) at (6.0,-2) {$|R'\rangle$};
    \draw[tensorred, thick] (L3) -- (D1);
    \draw[tensorred, thick] (D1) -- (D2);
    \draw[tensorred, thick] (D2) -- (D3);
    \draw[tensorred, thick] (D3) -- (D4);
\end{tikzpicture}
\caption{三类约束的 MPS 表示：蓝色为行约束（恰好一个），绿色为列约束（恰好一个），红色为对角线约束（至多一个）。键维度均为 2。}
\end{figure}

%========================================
\section{完整张量网络的构造}
%========================================

\subsection{多个约束在同一格点上的"汇聚"}

每个格点 $(i,j)$ 同时被\textbf{四条}约束链穿过：
\begin{enumerate}
    \item 第 $i$ 行的行约束（水平方向）
    \item 第 $j$ 列的列约束（垂直方向）
    \item 经过 $(i,j)$ 的 "$/$" 对角线约束
    \item 经过 $(i,j)$ 的 "$\backslash$" 对角线约束
\end{enumerate}

在每个格点上，这四条链共享同一个物理变量 $n_{ij}$。通过对 $n_{ij}$ 求和，我们把四条链"融合"成一个\textbf{格点张量}：

\begin{equation}
\boxed{
\mathcal{T}_{\alpha_L \alpha_R,\; \beta_U \beta_D,\; \gamma_{\text{in}} \gamma_{\text{out}},\; \delta_{\text{in}} \delta_{\text{out}}} = \sum_{n=0}^{1} A^{[n]}_{\alpha_L \alpha_R} \cdot A^{[n]}_{\beta_U \beta_D} \cdot B^{[n]}_{\gamma_{\text{in}} \gamma_{\text{out}}} \cdot B^{[n]}_{\delta_{\text{in}} \delta_{\text{out}}}
}
\end{equation}

其中：
\begin{itemize}
    \item $(\alpha_L, \alpha_R)$：水平键指标（行约束），维度 2
    \item $(\beta_U, \beta_D)$：垂直键指标（列约束），维度 2
    \item $(\gamma_{\text{in}}, \gamma_{\text{out}})$："$/$" 对角线键指标，维度 2
    \item $(\delta_{\text{in}}, \delta_{\text{out}})$："$\backslash$" 对角线键指标，维度 2
\end{itemize}

这个格点张量共有 $2^8 = 256$ 个元素，但大多数为零。

\subsection{显式计算格点张量}

由于 $n$ 只取 0 或 1，我们可以把求和写开：
\begin{align}
\mathcal{T} &= \underbrace{A^{[0]}_{\alpha_L \alpha_R} \cdot A^{[0]}_{\beta_U \beta_D} \cdot B^{[0]}_{\gamma\gamma'} \cdot B^{[0]}_{\delta\delta'}}_{n=0} + \underbrace{A^{[1]}_{\alpha_L \alpha_R} \cdot A^{[1]}_{\beta_U \beta_D} \cdot B^{[1]}_{\gamma\gamma'} \cdot B^{[1]}_{\delta\delta'}}_{n=1}
\end{align}

利用 $A^{[0]} = I$（单位矩阵）和 $B^{[0]} = I$，$n = 0$ 项为：
\[
n=0: \quad \delta_{\alpha_L \alpha_R} \cdot \delta_{\beta_U \beta_D} \cdot \delta_{\gamma\gamma'} \cdot \delta_{\delta\delta'}
\]
即所有键指标"穿过"——不放皇后时，各约束链的状态不改变。

$n = 1$ 项需要 $A^{[1]}$ 和 $B^{[1]}$ 都非零，即所有四个约束链都允许在此放皇后。由于 $A^{[1]}_{01} = 1$（其余为零）和 $B^{[1]}_{01} = 1$（其余为零），$n=1$ 项为：
\[
n=1: \quad \delta_{\alpha_L,0}\delta_{\alpha_R,1} \cdot \delta_{\beta_U,0}\delta_{\beta_D,1} \cdot \delta_{\gamma,0}\delta_{\gamma',1} \cdot \delta_{\delta,0}\delta_{\delta',1}
\]
即\textbf{所有四个键必须同时从状态 0 跳到状态 1}——这正是"在此放皇后"的张量网络表达。

\textbf{物理含义}：格点张量 $\mathcal{T}$ 实现了一个"逻辑开关"——
\begin{itemize}
    \item 不放皇后（$n=0$）：所有约束链的状态原样传递。
    \item 放皇后（$n=1$）：所有四个约束链必须从"未见皇后"状态（0）跳到"已见皇后"状态（1）。如果任何一条链已经处于状态 1（该行/列/对角线已有皇后），则这条路径的权重为 0——自动排除冲突。
\end{itemize}

\subsection{网络拓扑结构}

将所有格点张量按棋盘排列并连接相应的键，得到完整的张量网络：

\begin{figure}[H]
\centering
\begin{tikzpicture}[scale=1.1]
    % Draw 4x4 grid of tensors
    \foreach \i in {1,...,4} {
        \foreach \j in {1,...,4} {
            \node[tensor, fill=blue!10, minimum size=0.5cm, inner sep=1pt, font=\tiny] (T\i\j) at (\j*1.4, -\i*1.4) {};
        }
    }

    % Horizontal bonds (row constraints) - blue
    \foreach \i in {1,...,4} {
        \node[fill=yellow!50, circle, inner sep=1pt, font=\tiny] (HL\i) at (0.5, -\i*1.4) {$L$};
        \draw[tensorblue, thick] (HL\i) -- (T\i1);
        \foreach \j [count=\jn from 2] in {1,2,3} {
            \draw[tensorblue, thick] (T\i\j) -- (T\i\jn);
        }
        \node[fill=yellow!50, circle, inner sep=1pt, font=\tiny] (HR\i) at (6.1, -\i*1.4) {$R$};
        \draw[tensorblue, thick] (T\i4) -- (HR\i);
    }

    % Vertical bonds (column constraints) - green
    \foreach \j in {1,...,4} {
        \node[fill=yellow!50, circle, inner sep=1pt, font=\tiny] (VU\j) at (\j*1.4, -0.5) {$L$};
        \draw[tensorgreen, thick] (VU\j) -- (T1\j);
        \foreach \i [count=\in from 2] in {1,2,3} {
            \draw[tensorgreen, thick] (T\i\j) -- (T\in\j);
        }
        \node[fill=yellow!50, circle, inner sep=1pt, font=\tiny] (VD\j) at (\j*1.4, -6.1) {$R$};
        \draw[tensorgreen, thick] (T4\j) -- (VD\j);
    }

    % "\" diagonal bonds - red (upper-left to lower-right)
    \draw[tensorred, thick, opacity=0.5] (T11) -- (T22);
    \draw[tensorred, thick, opacity=0.5] (T22) -- (T33);
    \draw[tensorred, thick, opacity=0.5] (T33) -- (T44);

    \draw[tensorred, thick, opacity=0.5] (T12) -- (T23);
    \draw[tensorred, thick, opacity=0.5] (T23) -- (T34);

    \draw[tensorred, thick, opacity=0.5] (T13) -- (T24);

    \draw[tensorred, thick, opacity=0.5] (T21) -- (T32);
    \draw[tensorred, thick, opacity=0.5] (T32) -- (T43);

    \draw[tensorred, thick, opacity=0.5] (T31) -- (T42);

    % "/" diagonal bonds - orange (upper-right to lower-left)
    \draw[tensororange, thick, opacity=0.5] (T14) -- (T23);
    \draw[tensororange, thick, opacity=0.5] (T23) -- (T32);
    \draw[tensororange, thick, opacity=0.5] (T32) -- (T41);

    \draw[tensororange, thick, opacity=0.5] (T13) -- (T22);
    \draw[tensororange, thick, opacity=0.5] (T22) -- (T31);

    \draw[tensororange, thick, opacity=0.5] (T24) -- (T33);
    \draw[tensororange, thick, opacity=0.5] (T33) -- (T42);

    \draw[tensororange, thick, opacity=0.5] (T12) -- (T21);

    \draw[tensororange, thick, opacity=0.5] (T34) -- (T43);

    % Legend
    \node[font=\small] at (3.2, -7.2) {
        \textcolor{tensorblue}{--- 行约束} \quad
        \textcolor{tensorgreen}{--- 列约束} \quad
        \textcolor{tensorred}{--- $\backslash$ 对角线} \quad
        \textcolor{tensororange}{--- $/$ 对角线}
    };
\end{tikzpicture}
\caption{$4 \times 4$ 棋盘上的完整张量网络。每个节点是一个格点张量 $\mathcal{T}$，四种颜色的连线分别对应四类约束。$N$ 皇后解的总数 $Z_N$ 就是这个网络的完全收缩值。}
\end{figure}

\subsection{边界条件}

网络的边界由各约束链的端点向量决定：
\begin{itemize}
    \item \textbf{行约束}（水平方向）：左端 $\langle L| = (1,0)$，右端 $|R\rangle = (0,1)^T$。（恰好一个皇后。）
    \item \textbf{列约束}（垂直方向）：上端 $\langle L| = (1,0)$，下端 $|R\rangle = (0,1)^T$。（恰好一个皇后。）
    \item \textbf{对角线约束}：两端分别为 $\langle L| = (1,0)$ 和 $|R'\rangle = (1,1)^T$。（至多一个皇后。）
\end{itemize}

\textbf{配分函数}（$N$ 皇后解的总数）等于这个完整张量网络的收缩值：
\begin{equation}
\boxed{Z_N = \text{完全收缩}\Big(\text{所有格点张量 } \mathcal{T} \text{ + 边界向量}\Big)}
\end{equation}

%========================================
\section{转移矩阵视角}
%========================================

上一节的张量网络是"从空间角度"构造的。本节从"逐行扫描"的角度出发，导出等价的\textbf{转移矩阵}表述——这也是实际计算中最常用的方法。

\subsection{逐行处理的思路}

既然每行恰好一个皇后，我们可以用变量 $q_i \in \{1, \ldots, N\}$ 表示第 $i$ 行皇后所在的列。$N$ 皇后问题等价于：

\begin{tcolorbox}[colback=blue!5, colframe=blue!50!black]
找出所有满足以下条件的 $(q_1, q_2, \ldots, q_N) \in \{1, \ldots, N\}^N$：
\begin{enumerate}
    \item \textbf{列不重复}：$q_i \neq q_j$，$\forall\, i \neq j$
    \item \textbf{对角线不冲突}：$|q_i - q_j| \neq |i - j|$，$\forall\, i \neq j$
\end{enumerate}
\end{tcolorbox}

条件 1 意味着 $(q_1, \ldots, q_N)$ 是 $\{1, \ldots, N\}$ 的一个\textbf{排列}。条件 2 是额外的对角线限制。

如果我们能构造一个转移矩阵 $T$，使得 $Z_N = \langle \text{final}| \, T^N \, |\text{initial}\rangle$，那么就可以用张量网络的方法来计算它。

\subsection{为什么最近邻转移矩阵不够}

一个诱人的想法是定义一个 $N \times N$ 的"最近邻"转移矩阵：
\begin{equation}
T^{\text{nn}}_{ab} = \begin{cases} 1, & \text{若 } a \neq b \text{ 且 } |a - b| \neq 1 \\ 0, & \text{否则} \end{cases}
\end{equation}
它只检查\textbf{相邻两行}之间的约束。

然而，这忽略了两个关键的\textbf{长程约束}：
\begin{enumerate}
    \item \textbf{列约束是全局的}：即使 $q_1 \neq q_2$ 且 $q_2 \neq q_3$，仍可能 $q_1 = q_3$（列重复）。
    \item \textbf{对角线约束跨越多行}：第 1 行和第 3 行的皇后可能在同一对角线上（$|q_1 - q_3| = 2$），但 $T^{\text{nn}}$ 不检查这种情况。
\end{enumerate}

\textbf{例子}：$N = 4$，序列 $(1, 3, 1, 3)$。$T^{\text{nn}}$ 会接受它（每对相邻行都不冲突），但它不是合法的 $N$ 皇后解（列 1 和列 3 各出现了两次）。

\subsection{完整转移矩阵：增广状态空间}

为了正确编码所有约束，我们需要\textbf{增广}转移矩阵的状态空间，使其携带足够的"记忆"来追踪所有历史信息。

\subsubsection{状态的定义}

在放完第 $m$ 行的皇后之后，为了判断第 $m+1$ 行的哪些列可用，我们需要知道：

\begin{enumerate}
    \item \textbf{列占据向量} $\mathbf{c} = (c_1, \ldots, c_N) \in \{0,1\}^N$：$c_j = 1$ 表示列 $j$ 已被某行的皇后占据。这个信息\textbf{永久保留}——一旦某列被占，后续所有行都不能再用。

    \item \textbf{右对角线威胁向量} $\mathbf{r} = (r_1, \ldots, r_N) \in \{0,1\}^N$：$r_j = 1$ 表示列 $j$ 正受到来自左上方（"$\backslash$" 方向）的对角线攻击。当我们移到下一行时，这个威胁向右平移一格：$r'_j = r_{j-1}$（因为对角线每下一行就向右偏一格）。

    \item \textbf{左对角线威胁向量} $\mathbf{l} = (l_1, \ldots, l_N) \in \{0,1\}^N$：$l_j = 1$ 表示列 $j$ 正受到来自右上方（"$/$" 方向）的对角线攻击。移到下一行时向左平移一格：$l'_j = l_{j+1}$。
\end{enumerate}

完整状态为三元组 $|\mathbf{s}\rangle = |\mathbf{c}, \mathbf{r}, \mathbf{l}\rangle$，状态空间维度为 $2^{3N}$。

\subsubsection{状态更新规则}

在第 $m+1$ 行放置皇后于列 $q$。首先检查\textbf{合法性}：
\begin{equation}
\text{合法条件：} \quad c_q = 0 \;\wedge\; r_q = 0 \;\wedge\; l_q = 0
\end{equation}
即列 $q$ 未被占据，也不受任何对角线威胁。

然后\textbf{更新状态}（每个分量 $j = 1, \ldots, N$）：
\begin{align}
c'_j &= c_j \;\lor\; \delta_{j,q} \label{eq:col_update} \\
r'_j &= \begin{cases} r_{j-1} \;\lor\; \delta_{j, q+1}, & j \geq 2 \\ \delta_{j, q+1}, & j = 1 \end{cases} \label{eq:rdiag_update} \\
l'_j &= \begin{cases} l_{j+1} \;\lor\; \delta_{j, q-1}, & j \leq N-1 \\ \delta_{j, q-1}, & j = N \end{cases} \label{eq:ldiag_update}
\end{align}

\textbf{解读}：
\begin{itemize}
    \item 公式 \eqref{eq:col_update}：列 $q$ 被新皇后占据。
    \item 公式 \eqref{eq:rdiag_update}：旧的右对角线威胁向右平移一格（$r_{j-1} \to r'_j$），同时新皇后在列 $q$ 产生新的右对角线威胁，在下一行影响列 $q+1$。
    \item 公式 \eqref{eq:ldiag_update}：旧的左对角线威胁向左平移一格（$l_{j+1} \to l'_j$），新皇后产生新的左对角线威胁影响列 $q-1$。
\end{itemize}

转移矩阵为：
\begin{equation}
T = \sum_{q=1}^{N} T_q, \qquad T_q |\mathbf{c}, \mathbf{r}, \mathbf{l}\rangle = \begin{cases} |\mathbf{c}', \mathbf{r}', \mathbf{l}'\rangle & \text{若 $q$ 合法} \\ 0 & \text{否则} \end{cases}
\end{equation}

\subsection{配分函数的转移矩阵表达}

\begin{equation}
\boxed{
Z_N = \langle \text{final}| \; T^N \; |\text{initial}\rangle
}
\end{equation}

\begin{itemize}
    \item \textbf{初态}：$|\text{initial}\rangle = |\mathbf{0}, \mathbf{0}, \mathbf{0}\rangle$（所有列空闲，无对角线威胁）。
    \item \textbf{末态}：$\langle \text{final}| = \langle \mathbf{1}| \otimes \langle \text{any}| \otimes \langle \text{any}|$，即所有列都被占据（$c_j = 1$，$\forall j$），对角线状态无要求。
\end{itemize}

\textbf{等价地}：$\langle \text{final}|$ 投影到 $c_j = 1, \forall j$ 的子空间。由于我们精确地放了 $N$ 个皇后在 $N$ 列中，最终必然所有列都被占据，所以这个约束自动满足。实际上可以简单取 $\langle \text{final}| = \sum_{\mathbf{r}, \mathbf{l}} \langle \mathbf{1}, \mathbf{r}, \mathbf{l}|$。

%========================================
\section{转移矩阵的 MPO 分解}
%========================================

$2^{3N}$ 维的转移矩阵看似巨大，但它具有\textbf{张量积结构}——这正是张量网络方法的切入点。

\subsection{状态空间的张量积分解}

完整状态可以分解为 $N$ 个局部空间的张量积：
\begin{equation}
|\mathbf{s}\rangle = |s_1\rangle \otimes |s_2\rangle \otimes \cdots \otimes |s_N\rangle, \qquad s_j = (c_j, r_j, l_j) \in \{0,1\}^3
\end{equation}

每个列 $j$ 对应一个 $2^3 = 8$ 维的局部希尔伯特空间。

\subsection{MPO 结构}

转移矩阵 $T = \sum_q T_q$ 可以表示为\textbf{矩阵乘积算符}（MPO）：
\begin{equation}
T_{s'_1 \cdots s'_N,\; s_1 \cdots s_N} = \sum_{\beta_0, \beta_1, \ldots, \beta_N} L_{\beta_0} \cdot W^{[1]}_{s'_1, s_1, \beta_0, \beta_1} \cdot W^{[2]}_{s'_2, s_2, \beta_1, \beta_2} \cdots W^{[N]}_{s'_N, s_N, \beta_{N-1}, \beta_N} \cdot R_{\beta_N}
\end{equation}

其中 $W^{[j]}$ 是列 $j$ 的\textbf{MPO 局部张量}，$\beta$ 是水平键指标（MPO 键维度），$L$ 和 $R$ 是边界向量。

\textbf{MPO 键指标 $\beta$ 的含义}：水平键需要传递两类信息：
\begin{enumerate}
    \item \textbf{皇后放置标记}：在扫描到列 $j$ 时，皇后是否已经放好？$p \in \{0, 1\}$。
    \item \textbf{右对角线传递}：列 $j-1$ 的旧右对角线状态 $r_{j-1}$，需要传给列 $j$ 用于更新 $r'_j$。值为 $\{0, 1\}$。
\end{enumerate}

因此 $\beta = (p, r_{\text{pass}}) \in \{0,1\}^2$，MPO 键维度为 $D_{\text{MPO}} = 4$。

\textbf{左对角线的处理}：左对角线威胁 $l'_j = l_{j+1}$ 需要从右边传递信息。有两种处理方式：
\begin{enumerate}
    \item 将转移矩阵分解为\textbf{两个 MPO 的乘积}：第一个（从左到右）处理列约束和右对角线，第二个（从右到左）处理左对角线。
    \item 直接读取纵向键上的 $l_{j+1}$（旧值），因为 $l_{j+1}$ 是上一行传下来的已知信息，存储在纵向键中。
\end{enumerate}

无论哪种方式，MPO 的总键维度都是\textbf{常数}（不随 $N$ 增长）。

\subsection{MPO 局部张量的显式形式}

列 $j$ 的 MPO 张量 $W^{[j]}$ 作用在：
\begin{itemize}
    \item 纵向输入 $s_j = (c_j, r_j, l_j)$（从上一行传来）
    \item 纵向输出 $s'_j = (c'_j, r'_j, l'_j)$（传给下一行）
    \item 水平输入 $\beta_{j-1} = (p, r_{\text{pass}})$
    \item 水平输出 $\beta_j = (p', r'_{\text{pass}})$
\end{itemize}

$W^{[j]}$ 的非零元素由以下逻辑决定：

\begin{tcolorbox}[colback=gray!5, colframe=gray!60, title=MPO 局部张量的逻辑]
\textbf{情况 1：不在列 $j$ 放皇后}（$p' = p$，皇后标记不变）
\begin{itemize}
    \item $c'_j = c_j$（列状态不变）
    \item $r'_j = r_{\text{pass}}$（接收左邻的右对角线信息）
    \item $l'_j = l_{j+1}$（从纵向键读取右邻上一行的左对角线信息）
    \item $r'_{\text{pass}} = r_j$（将本列的旧右对角线传给右邻）
\end{itemize}

\textbf{情况 2：在列 $j$ 放皇后}（$p = 0 \to p' = 1$，且 $c_j = 0, r_q = 0, l_q = 0$）
\begin{itemize}
    \item $c'_j = 1$
    \item $r'_j = r_{\text{pass}}$（旧威胁仍在）
    \item $l'_j = l_{j+1}$（旧威胁仍在）
    \item $r'_{\text{pass}} = r_j \;\lor\; 1 = 1$（新的右对角线威胁产生）
    \item 同时需要向左传递新的左对角线威胁
\end{itemize}
\end{tcolorbox}

\textbf{关键性质}：MPO 键维度为 $O(1)$——不随 $N$ 增长。这意味着转移矩阵虽然是 $2^{3N} \times 2^{3N}$ 的巨大矩阵，但具有非常紧凑的 MPO 表示。

%========================================
\section{张量网络的收缩}
%========================================

\subsection{精确收缩：MPS-MPO 方法}

计算 $Z_N = \langle \text{final}| \, T^N \, |\text{initial}\rangle$ 的过程可以理解为：

\begin{enumerate}
    \item 初始化一个 MPS $|\psi_0\rangle = |\text{initial}\rangle = |\mathbf{0}\rangle$（键维度 $\chi = 1$）。
    \item 逐行施加转移矩阵 MPO：$|\psi_m\rangle = T \, |\psi_{m-1}\rangle$。
    \item 最后与末态内积：$Z_N = \langle \text{final}|\psi_N\rangle$。
\end{enumerate}

\begin{figure}[H]
\centering
\begin{tikzpicture}[scale=0.85]
    % MPS at top
    \node[font=\small] at (-1.5, 4) {$|\psi_0\rangle$};
    \foreach \j in {1,...,5} {
        \node[tensor, fill=green!20, minimum size=0.5cm] (M0\j) at (\j*1.3, 4) {};
    }
    \foreach \j [count=\jn from 2] in {1,...,4} {
        \draw[thick] (M0\j) -- (M0\jn);
    }

    % MPO layer 1
    \node[font=\small] at (-1.5, 2.8) {$T$（第1行）};
    \foreach \j in {1,...,5} {
        \node[tensor, fill=orange!20, minimum size=0.5cm] (T1\j) at (\j*1.3, 2.8) {};
        \draw[thick] (M0\j) -- (T1\j);
    }
    \foreach \j [count=\jn from 2] in {1,...,4} {
        \draw[thick] (T1\j) -- (T1\jn);
    }

    % MPO layer 2
    \node[font=\small] at (-1.5, 1.6) {$T$（第2行）};
    \foreach \j in {1,...,5} {
        \node[tensor, fill=orange!20, minimum size=0.5cm] (T2\j) at (\j*1.3, 1.6) {};
        \draw[thick] (T1\j) -- (T2\j);
    }
    \foreach \j [count=\jn from 2] in {1,...,4} {
        \draw[thick] (T2\j) -- (T2\jn);
    }

    % dots
    \node at (3.9, 0.7) {$\vdots$};

    % MPO layer N
    \node[font=\small] at (-1.5, -0.2) {$T$（第$N$行）};
    \foreach \j in {1,...,5} {
        \node[tensor, fill=orange!20, minimum size=0.5cm] (TN\j) at (\j*1.3, -0.2) {};
    }
    \foreach \j [count=\jn from 2] in {1,...,4} {
        \draw[thick] (TN\j) -- (TN\jn);
    }

    % Final MPS
    \node[font=\small] at (-1.5, -1.4) {$\langle\text{final}|$};
    \foreach \j in {1,...,5} {
        \node[tensor, fill=green!20, minimum size=0.5cm] (MF\j) at (\j*1.3, -1.4) {};
        \draw[thick] (TN\j) -- (MF\j);
    }
    \foreach \j [count=\jn from 2] in {1,...,4} {
        \draw[thick] (MF\j) -- (MF\jn);
    }
\end{tikzpicture}
\caption{$Z_N$ 的 MPS-MPO 收缩结构：初始 MPS（绿色）经过 $N$ 层 MPO（橙色）后与末态 MPS 收缩。}
\end{figure}

\subsection{键维度的增长}

每次 MPS-MPO 相乘，MPS 的键维度原则上会乘以 MPO 的键维度。设 MPS 键维度为 $\chi_m$，MPO 键维度为 $D_{\text{MPO}}$，则：
\begin{equation}
\chi_{m+1} \leq \chi_m \times D_{\text{MPO}} \times d
\end{equation}
其中 $d = 8$ 是局部希尔伯特空间维度。

对于精确计算，$\chi$ 的增长由\textbf{可达状态数}决定。放完 $m$ 行后：
\begin{itemize}
    \item 恰好 $m$ 列被占据（列占据向量中恰好 $m$ 个 1）：$\binom{N}{m}$ 种选择。
    \item 对角线威胁模式受限于前 $m$ 行的皇后位置。
\end{itemize}

精确的 MPS 键维度在 $m \approx N/2$ 时达到峰值，量级为 $O\!\left(\binom{N}{N/2}\right) \sim O(2^N/\sqrt{N})$——\textbf{仍然是指数增长的}。这是意料之中的：\#$N$-Queens 是 \#P-完全的，不可能用多项式资源精确求解。

\subsection{近似收缩：MPS 截断}

对于大 $N$，我们可以在每次 MPS-MPO 相乘后，通过\textbf{SVD 截断}将 MPS 键维度限制在 $\chi_{\max}$ 以内：

\begin{tcolorbox}[colback=yellow!5, colframe=orange!60!black, title=\textbf{近似收缩算法}]
\begin{enumerate}[leftmargin=*]
    \item $|\psi_0\rangle = |\mathbf{0}\rangle$，$\chi = 1$
    \item \textbf{for} $m = 1, 2, \ldots, N$:
    \begin{enumerate}
        \item $|\tilde{\psi}_m\rangle = T \, |\psi_{m-1}\rangle$（MPS-MPO 乘法，$\chi \to \chi \cdot D_{\text{MPO}}$）
        \item \textbf{if} $\chi > \chi_{\max}$：对 $|\tilde{\psi}_m\rangle$ 做 SVD 截断，保留前 $\chi_{\max}$ 个奇异值
        \item $|\psi_m\rangle = |\tilde{\psi}_m\rangle$（截断后的 MPS）
    \end{enumerate}
    \item $Z_N \approx \langle \text{final}|\psi_N\rangle$
\end{enumerate}
\end{tcolorbox}

截断引入的误差取决于被丢弃的奇异值的大小。$\chi_{\max}$ 越大，结果越精确。

\textbf{注意}：与物理系统（如 Ising 模型）不同，$N$ 皇后问题的"纠缠"结构较为特殊——长程约束（特别是列约束"所有列不重复"）意味着左半和右半棋盘之间存在强关联，这使得 MPS 截断的效率可能不如物理系统中那么高。

%========================================
\section{具体计算示例：$N = 4$}
%========================================

让我们对 $N = 4$ 的情况做完整的转移矩阵计算，追踪每一步的状态变化。

\subsection{$4 \times 4$ 棋盘的两个解}

$N = 4$ 有且仅有 2 个解：

\begin{figure}[H]
\centering
\begin{tikzpicture}[scale=0.8]
    % Solution 1: (2,4,1,3)
    \node[font=\bfseries] at (2.5, 1) {解 1：$(2,4,1,3)$};
    \foreach \i in {1,...,4} {
        \foreach \j in {1,...,4} {
            \pgfmathtruncatemacro{\shade}{mod(\i+\j,2)*20}
            \node[draw, minimum size=0.9cm, fill=gray!\shade] (A\i\j) at (\j, -\i) {};
        }
    }
    \node[queen] at (2, -1) {Q};
    \node[queen] at (4, -2) {Q};
    \node[queen] at (1, -3) {Q};
    \node[queen] at (3, -4) {Q};

    % Solution 2: (3,1,4,2)
    \begin{scope}[xshift=6.5cm]
    \node[font=\bfseries] at (2.5, 1) {解 2：$(3,1,4,2)$};
    \foreach \i in {1,...,4} {
        \foreach \j in {1,...,4} {
            \pgfmathtruncatemacro{\shade}{mod(\i+\j,2)*20}
            \node[draw, minimum size=0.9cm, fill=gray!\shade] (B\i\j) at (\j, -\i) {};
        }
    }
    \node[queen] at (3, -1) {Q};
    \node[queen] at (1, -2) {Q};
    \node[queen] at (4, -3) {Q};
    \node[queen] at (2, -4) {Q};
    \end{scope}
\end{tikzpicture}
\caption{$N=4$ 的两个解：$(q_1,q_2,q_3,q_4) = (2,4,1,3)$ 和 $(3,1,4,2)$。}
\end{figure}

\subsection{逐行追踪状态}

以解 1：$(2, 4, 1, 3)$ 为例，追踪转移矩阵的状态 $(\mathbf{c}, \mathbf{r}, \mathbf{l})$。

\subsubsection*{初态（放置第 1 行之前）}

\[
\mathbf{s}_0 = \Big(\underbrace{0,0,0,0}_{\mathbf{c}},\; \underbrace{0,0,0,0}_{\mathbf{r}},\; \underbrace{0,0,0,0}_{\mathbf{l}}\Big)
\]
棋盘完全空白，所有列都可用，无对角线威胁。

\subsubsection*{第 1 行：放皇后于列 2}

检查合法性：$c_2 = 0$, $r_2 = 0$, $l_2 = 0$ \checkmark

更新：
\begin{align*}
\mathbf{c}' &= (0, \mathbf{1}, 0, 0) \quad\text{（列 2 被占据）} \\
\mathbf{r}' &: \quad r'_j = r_{j-1} \;\lor\; \delta_{j,3} \\
&\quad r'_1 = 0,\; r'_2 = r_1 = 0,\; r'_3 = r_2 \lor 1 = \mathbf{1},\; r'_4 = r_3 = 0 \\
&= (0, 0, \mathbf{1}, 0) \quad\text{（列 3 受右对角线威胁）} \\
\mathbf{l}' &: \quad l'_j = l_{j+1} \;\lor\; \delta_{j,1} \\
&\quad l'_1 = l_2 \lor 1 = \mathbf{1},\; l'_2 = l_3 = 0,\; l'_3 = l_4 = 0,\; l'_4 = 0 \\
&= (\mathbf{1}, 0, 0, 0) \quad\text{（列 1 受左对角线威胁）}
\end{align*}

\[
\mathbf{s}_1 = (0,1,0,0;\; 0,0,1,0;\; 1,0,0,0)
\]

\textbf{可用列}（$c_j = 0$ 且 $r_j = 0$ 且 $l_j = 0$）：
\begin{center}
\begin{tabular}{c|cccc}
列 $j$ & 1 & 2 & 3 & 4 \\ \hline
$c_j$ & 0 & 1 & 0 & 0 \\
$r_j$ & 0 & 0 & 1 & 0 \\
$l_j$ & 1 & 0 & 0 & 0 \\
可用？ & $\times$ & $\times$ & $\times$ & $\checkmark$
\end{tabular}
\end{center}
只有列 4 可用！

\subsubsection*{第 2 行：放皇后于列 4}

检查合法性：$c_4 = 0$, $r_4 = 0$, $l_4 = 0$ \checkmark

更新：
\begin{align*}
\mathbf{c}' &= (0, 1, 0, \mathbf{1}) \\
\mathbf{r}' &: \quad r'_1 = 0,\; r'_2 = r_1 = 0,\; r'_3 = r_2 = 0,\; r'_4 = r_3 \lor 0 = \mathbf{1} \\
&\quad\text{（注意：$\delta_{j,5} = 0$ 因为列 5 不存在；旧威胁 $r_3 = 1$ 移到 $r'_4$）} \\
&= (0, 0, 0, \mathbf{1}) \\
\mathbf{l}' &: \quad l'_1 = l_2 = 0,\; l'_2 = l_3 = 0,\; l'_3 = l_4 \lor 1 = \mathbf{1},\; l'_4 = 0 \\
&= (0, 0, \mathbf{1}, 0)
\end{align*}

\[
\mathbf{s}_2 = (0,1,0,1;\; 0,0,0,1;\; 0,0,1,0)
\]

可用列：
\begin{center}
\begin{tabular}{c|cccc}
列 $j$ & 1 & 2 & 3 & 4 \\ \hline
$c_j$ & 0 & 1 & 0 & 1 \\
$r_j$ & 0 & 0 & 0 & 1 \\
$l_j$ & 0 & 0 & 1 & 0 \\
可用？ & $\checkmark$ & $\times$ & $\times$ & $\times$
\end{tabular}
\end{center}
只有列 1 可用！

\subsubsection*{第 3 行：放皇后于列 1}

检查合法性：$c_1 = 0$, $r_1 = 0$, $l_1 = 0$ \checkmark

更新：
\begin{align*}
\mathbf{c}' &= (\mathbf{1}, 1, 0, 1) \\
\mathbf{r}' &: \quad r'_1 = 0,\; r'_2 = r_1 \lor 1 = \mathbf{1},\; r'_3 = r_2 = 0,\; r'_4 = r_3 = 0 \\
&\quad\text{（旧 $r_4 = 1$ 移出棋盘；新威胁影响列 2）} \\
&= (0, \mathbf{1}, 0, 0) \\
\mathbf{l}' &: \quad l'_1 = l_2 = 0,\; l'_2 = l_3 = 1,\; l'_3 = l_4 = 0,\; l'_4 = 0 \\
&\quad\text{（新威胁 $\delta_{j,0}$：列 0 不存在，无新左对角线威胁）} \\
&= (0, \mathbf{1}, 0, 0)
\end{align*}

等一下，让我重新检查 $l'_2$。旧的 $\mathbf{l} = (0, 0, 1, 0)$。$l'_j = l_{j+1} \lor \delta_{j, q-1}$，$q = 1$，所以 $\delta_{j, 0} = 0$ 对所有 $j$。
\begin{align*}
l'_1 &= l_2 = 0 \\
l'_2 &= l_3 = 1 \\
l'_3 &= l_4 = 0 \\
l'_4 &= 0
\end{align*}
所以 $\mathbf{l}' = (0, 1, 0, 0)$。这里 $l'_2 = 1$ 来自第 2 行列 4 皇后的左对角线威胁（经过两步平移：$l_4 \to l_3 \to l_2$）。

\[
\mathbf{s}_3 = (1,1,0,1;\; 0,1,0,0;\; 0,1,0,0)
\]

可用列：
\begin{center}
\begin{tabular}{c|cccc}
列 $j$ & 1 & 2 & 3 & 4 \\ \hline
$c_j$ & 1 & 1 & 0 & 1 \\
$r_j$ & 0 & 1 & 0 & 0 \\
$l_j$ & 0 & 1 & 0 & 0 \\
可用？ & $\times$ & $\times$ & $\checkmark$ & $\times$
\end{tabular}
\end{center}
只有列 3 可用！

\subsubsection*{第 4 行：放皇后于列 3}

$c_3 = 0$, $r_3 = 0$, $l_3 = 0$ \checkmark

\[
\mathbf{s}_4 = (1,1,1,1;\; \cdot,\cdot,\cdot,\cdot;\; \cdot,\cdot,\cdot,\cdot)
\]

所有列都被占据——这是一个合法的 $N$ 皇后解！

\subsection{状态演化的张量网络图}

我们可以把上述过程画成一个张量网络图：

\begin{figure}[H]
\centering
\begin{tikzpicture}[scale=1.0]
    % Labels
    \node[font=\small] at (-0.5, 0.5) {列};
    \node[font=\small] at (-2.2, 0.5) {行};
    \foreach \j in {1,...,4} {
        \node[font=\small] at (\j*1.5, 0.5) {$j=\j$};
    }

    % Row 1
    \node[font=\small] at (-2.2, 0) {$i=1$};
    \foreach \j in {1,...,4} {
        \pgfmathtruncatemacro{\isqueen}{(\j==2) ? 1 : 0}
        \ifnum\isqueen=1
            \node[tensor, fill=red!30, minimum size=0.7cm] (T1\j) at (\j*1.5, 0) {\tiny Q};
        \else
            \node[tensor, fill=blue!10, minimum size=0.7cm] (T1\j) at (\j*1.5, 0) {};
        \fi
    }

    % Row 2
    \node[font=\small] at (-2.2, -1.5) {$i=2$};
    \foreach \j in {1,...,4} {
        \pgfmathtruncatemacro{\isqueen}{(\j==4) ? 1 : 0}
        \ifnum\isqueen=1
            \node[tensor, fill=red!30, minimum size=0.7cm] (T2\j) at (\j*1.5, -1.5) {\tiny Q};
        \else
            \node[tensor, fill=blue!10, minimum size=0.7cm] (T2\j) at (\j*1.5, -1.5) {};
        \fi
    }

    % Row 3
    \node[font=\small] at (-2.2, -3) {$i=3$};
    \foreach \j in {1,...,4} {
        \pgfmathtruncatemacro{\isqueen}{(\j==1) ? 1 : 0}
        \ifnum\isqueen=1
            \node[tensor, fill=red!30, minimum size=0.7cm] (T3\j) at (\j*1.5, -3) {\tiny Q};
        \else
            \node[tensor, fill=blue!10, minimum size=0.7cm] (T3\j) at (\j*1.5, -3) {};
        \fi
    }

    % Row 4
    \node[font=\small] at (-2.2, -4.5) {$i=4$};
    \foreach \j in {1,...,4} {
        \pgfmathtruncatemacro{\isqueen}{(\j==3) ? 1 : 0}
        \ifnum\isqueen=1
            \node[tensor, fill=red!30, minimum size=0.7cm] (T4\j) at (\j*1.5, -4.5) {\tiny Q};
        \else
            \node[tensor, fill=blue!10, minimum size=0.7cm] (T4\j) at (\j*1.5, -4.5) {};
        \fi
    }

    % Vertical bonds
    \foreach \j in {1,...,4} {
        \draw[tensorgreen, thick] (T1\j) -- (T2\j);
        \draw[tensorgreen, thick] (T2\j) -- (T3\j);
        \draw[tensorgreen, thick] (T3\j) -- (T4\j);
    }

    % Horizontal bonds
    \foreach \i in {1,...,4} {
        \draw[tensorblue, thick] (T\i1) -- (T\i2);
        \draw[tensorblue, thick] (T\i2) -- (T\i3);
        \draw[tensorblue, thick] (T\i3) -- (T\i4);
    }

    % State annotations
    \node[font=\tiny, right] at (7.2, 0) {$\mathbf{c}=(0,1,0,0)$};
    \node[font=\tiny, right] at (7.2, -1.5) {$\mathbf{c}=(0,1,0,1)$};
    \node[font=\tiny, right] at (7.2, -3) {$\mathbf{c}=(1,1,0,1)$};
    \node[font=\tiny, right] at (7.2, -4.5) {$\mathbf{c}=(1,1,1,1)$};
\end{tikzpicture}
\caption{解 $(2,4,1,3)$ 对应的张量网络状态演化。红色标记的格点放置了皇后，绿色纵向键携带列/对角线约束信息，蓝色横向键携带行约束信息。右侧标注了各行之后的列占据状态。}
\end{figure}

\subsection{计数验证}

对 $N = 4$，解 1 的路径 $(2,4,1,3)$ 中：
\begin{itemize}
    \item 第 1 步有 4 个可选列（$q_1 \in \{1,2,3,4\}$），选 $q_1 = 2$。
    \item 选 $q_1 = 2$ 后只有 $q_2 = 4$ 可用。
    \item 之后 $q_3 = 1$、$q_4 = 3$ 均唯一确定。
\end{itemize}

类似分析，选 $q_1 = 3$ 会导出解 2：$(3,1,4,2)$。选 $q_1 = 1$ 或 $q_1 = 4$ 最终会发现无列可用（死胡同）。所以 $Z_4 = 2$。

\textbf{转移矩阵验证}：由于 $N = 4$ 的状态空间为 $2^{12} = 4096$ 维，直接构造 $T$ 并计算 $T^4$ 是可行的（虽然不必手算），其结果确实给出 $Z_4 = 2$。

%========================================
\section{键维度与计算复杂度分析}
%========================================

\subsection{精确 MPS 键维度的增长}

在逐行收缩过程中，MPS 在第 $m$ 行和第 $m+1$ 行之间的键维度 $\chi_m$ 等于可达状态数。下表列出了 $N = 8$ 的情况：

\begin{center}
\renewcommand{\arraystretch}{1.2}
\begin{tabular}{c|ccccccccc}
\hline
行 $m$ & 0 & 1 & 2 & 3 & 4 & 5 & 6 & 7 & 8 \\
\hline
$\chi_m$（精确） & 1 & 8 & 44 & 172 & 436 & 568 & 316 & 76 & 1 \\
\hline
\end{tabular}
\end{center}

\textbf{观察}：
\begin{itemize}
    \item $\chi_0 = 1$（初始状态唯一）。
    \item $\chi_1 = 8 = N$（第 1 行皇后可放 8 个位置中的任一个）。
    \item $\chi$ 在 $m \approx N/2$ 处达到峰值。
    \item $\chi_N = 1$（若只跟踪合法构型；或 $\chi_N = Z_N$ 若跟踪所有路径）。
    \item 精确 MPS 键维度远小于 $2^{3N}$（$N=8$ 时为 $16777216$），但仍随 $N$ 指数增长。
\end{itemize}

\subsection{与其他方法的复杂度对比}

\begin{center}
\renewcommand{\arraystretch}{1.3}
\begin{tabular}{lcc}
\hline
\textbf{方法} & \textbf{时间复杂度} & \textbf{空间复杂度} \\
\hline
暴力枚举 & $O(N^N)$ & $O(N)$ \\
回溯搜索 & $O(N!)$（最坏） & $O(N)$ \\
转移矩阵（精确） & $O(N \cdot 2^{3N})$ & $O(2^{3N})$ \\
MPS-MPO（精确） & $O(N \cdot \chi_{\max}^2 \cdot D^2)$ & $O(N \cdot \chi_{\max})$ \\
MPS-MPO（截断到 $\chi$） & $O(N^2 \cdot \chi^3)$ & $O(N \cdot \chi)$ \\
\hline
\end{tabular}
\end{center}

\textbf{MPS-MPO 方法的优势}：
\begin{itemize}
    \item 精确版本的 $\chi_{\max}$ 虽然指数增长，但通常远小于 $2^{3N}$（因为大量状态不可达）。
    \item 截断版本提供了\textbf{可调精度的近似}——增大 $\chi$ 可以系统地提高精度。
    \item 张量网络框架统一了精确计数和近似计数，且易于推广到变体问题。
\end{itemize}

%========================================
\section{推广：皇后格点气体模型}
%========================================

$N$ 皇后问题要求恰好放置 $N$ 个皇后（每行每列各一个）。一个自然的推广是\textbf{皇后格点气体}（queen lattice gas）：在 $N \times N$ 棋盘上放置\textbf{任意数量}的互不攻击的皇后，每个皇后携带逸度 $z$。

\subsection{模型定义}

配分函数为：
\begin{equation}
\boxed{
\mathcal{Z}(z) = \sum_{k=0}^{N} z^k \cdot Q(N, k)
}
\end{equation}
其中 $Q(N, k)$ 是在 $N \times N$ 棋盘上放置 $k$ 个互不攻击皇后的方案数。特别地，$Q(N, N) = Z_N$ 就是 $N$ 皇后数。

等价地，使用格点变量：
\begin{equation}
\mathcal{Z}(z) = \sum_{\{n_{ij}\}} z^{\sum_{ij} n_{ij}} \prod_{\text{约束}} \delta(\text{无冲突})
\end{equation}

\textbf{物理图像}：这是一个\textbf{硬核格点气体}——皇后之间的互不攻击约束扮演了无穷大排斥相互作用（硬核势）的角色。逸度 $z$ 控制系统的密度：$z \to 0$ 时棋盘几乎为空，$z \to \infty$ 时系统趋于最密堆积。

\subsection{张量网络表示}

格点气体的张量网络与 $N$ 皇后问题几乎完全相同，唯一的区别是：

\begin{enumerate}
    \item \textbf{行约束放宽}：从"恰好一个"变为"至多一个"。MPS 的右边界向量从 $|R\rangle = (0,1)^T$ 改为 $|R'\rangle = (1,1)^T$。
    \item \textbf{列约束放宽}：同样从"恰好一个"变为"至多一个"。
    \item \textbf{逸度权重}：每放一个皇后，贡献一个因子 $z$。这可以通过修改局部张量实现：
\end{enumerate}

\begin{equation}
A_z^{[0]} = \begin{pmatrix} 1 & 0 \\ 0 & 1 \end{pmatrix}, \qquad
A_z^{[1]} = \begin{pmatrix} 0 & \sqrt[4]{z} \\ 0 & 0 \end{pmatrix}
\end{equation}

这里将 $z$ 的四分之一次方分配给每个方向的约束链（行、列、两个对角线各贡献 $z^{1/4}$），使得每个皇后的总权重为 $(z^{1/4})^4 = z$。

格点张量变为：
\begin{equation}
\mathcal{T}_z = \sum_{n=0,1} A_z^{[n]}_{\text{row}} \cdot A_z^{[n]}_{\text{col}} \cdot B_z^{[n]}_{\text{diag}/} \cdot B_z^{[n]}_{\text{diag}\backslash}
\end{equation}

\subsection{相变行为}

皇后格点气体在热力学极限 $N \to \infty$ 下展现出丰富的相变行为。

\subsubsection*{低密度相（$z < z_c$）}

当逸度 $z$ 较小时，皇后密度低，约束效应不强，系统类似于普通的格点气体。皇后在棋盘上近似均匀分布，没有长程序。

\subsubsection*{高密度相（$z > z_c$）}

当 $z$ 超过某个临界值 $z_c$ 时，皇后密度高到一定程度，系统被迫形成有序结构以满足互不攻击的约束。这种现象类似于\textbf{熵驱动的有序化}——硬核排斥本身就能导致相变。

\subsubsection*{张量网络方法的优势}

蒙特卡洛模拟是研究皇后格点气体相变的主要数值工具，但在高密度区域面临\textbf{严重的遍历性问题}——合法构型之间的跃迁极其困难。张量网络方法没有这个问题，因为它直接计算配分函数（或其近似），不依赖于构型空间的遍历性。

具体来说，MPS-MPO 方法可以：
\begin{itemize}
    \item 直接计算 $\mathcal{Z}(z)$ 作为 $z$ 的函数
    \item 通过 $\mathcal{Z}(z)$ 的解析性质（李-杨零点分布）定位相变点
    \item 在临界点附近通过增大 $\chi$ 系统地提高精度
    \item 计算纠缠熵等张量网络特有的诊断量，帮助理解相变的性质
\end{itemize}

%========================================
\section{讨论与展望}
%========================================

\subsection{方法总结}

\begin{tcolorbox}[colback=blue!5, colframe=blue!60!black, title=\textbf{$N$ 皇后问题的张量网络解法：完整流程}]
\begin{enumerate}
    \item \textbf{变量与约束}：在 $N \times N$ 格点上定义二值变量 $n_{ij}$，识别行/列/对角线约束。
    \item \textbf{约束编码}：每个约束用键维度 $D=2$ 的 MPS 表示（"恰好一个"或"至多一个"）。
    \item \textbf{格点张量}：在每个格点对物理指标求和，融合四条约束链为一个 8-键张量。
    \item \textbf{张量网络}：所有格点张量按棋盘排列并连接，形成具有水平、垂直、对角线键的二维张量网络。
    \item \textbf{转移矩阵}：逐行收缩，转移矩阵追踪列占据和对角线威胁状态，具有 MPO 结构。
    \item \textbf{收缩}：通过 MPS-MPO 乘法逐行收缩网络。精确收缩给出精确解；截断收缩给出可控近似。
    \item \textbf{结果}：$Z_N$ = 网络的收缩值 = $N$ 皇后解的总数。
\end{enumerate}
\end{tcolorbox}

\subsection{张量网络 vs. 传统算法}

\begin{center}
\renewcommand{\arraystretch}{1.3}
\begin{tabular}{l|ccc}
\hline
& \textbf{回溯搜索} & \textbf{转移矩阵} & \textbf{张量网络} \\
\hline
思想 & 逐行试错 & 状态空间枚举 & 约束分解与收缩 \\
精确性 & 精确 & 精确 & 精确或近似 \\
时间 & $O(N!)$ & $O(N \cdot 2^{3N})$ & $O(N \cdot \chi^3)$ \\
可推广性 & 低 & 中 & 高 \\
物理联系 & 无 & 量子力学类比 & 统计力学类比 \\
\hline
\end{tabular}
\end{center}

张量网络方法的独特价值不在于比回溯搜索更快地精确计数（对于精确计数，回溯搜索通常更实用），而在于：
\begin{enumerate}
    \item 提供了一种\textbf{系统化的近似框架}，可以通过增大 $\chi$ 逐步提高精度。
    \item 自然地与\textbf{统计力学}和\textbf{量子信息}的语言对接，使组合问题的物理图像更加清晰。
    \item 容易推广到\textbf{变体问题}（格点气体、加权计数、不同棋盘几何等）。
    \item 通过纠缠结构分析，揭示问题的\textbf{内在复杂度}——MPS 键维度的增长速率反映了约束的非局域程度。
\end{enumerate}

\subsection{进一步的方向}

\begin{itemize}
    \item \textbf{TRG/HOTRG 方法}：直接收缩二维张量网络（不经过 MPS-MPO），利用粗粒化减小网络尺寸。
    \item \textbf{CTMRG 方法}：用角转移矩阵近似环境，特别适合研究热力学极限下的格点气体。
    \item \textbf{PEPS 表示}：将解的集合表示为投影纠缠对态，利用其二维结构。
    \item \textbf{量子计算}：$N$ 皇后问题的张量网络结构可以映射到量子电路，利用量子计算机加速。
\end{itemize}

\vspace{0.5cm}

\begin{thebibliography}{99}
\bibitem{Orus2014} R. Or\'{u}s, ``A practical introduction to tensor networks: Matrix product states and projected entangled pair states,'' \textit{Ann. Phys.} \textbf{349}, 117 (2014).
\bibitem{Schollwock2011} U. Schollw\"{o}ck, ``The density-matrix renormalization group in the age of matrix product states,'' \textit{Ann. Phys.} \textbf{326}, 96 (2011).
\bibitem{Murg2005} V. Murg, F. Verstraete, and J. I. Cirac, ``Efficient evaluation of partition functions of frustrated and inhomogeneous spin systems,'' \textit{Phys. Rev. Lett.} \textbf{95}, 057206 (2005).
\bibitem{Kourtis2019} S. Kourtis, C. Chamon, E. R. Mucciolo, and A. E. Ruckenstein, ``Fast counting with tensor networks,'' \textit{SciPost Phys.} \textbf{7}, 060 (2019).
\bibitem{GarciaMartin2021} D. Garc\'{i}a-Mart\'{i}n and G. Sierra, ``Quantum computation of the classically intractable permanent using tensor networks,'' \textit{Phys. Rev. A} \textbf{104}, 032417 (2021).
\bibitem{Zhao2010} J. Zhao, ``Queens on board,'' OEIS A000170.
\end{thebibliography}

\end{document}
